\section{Discussion}
\label{sec:disc}

\paragraph{Why multiplicative?}
The cleanest distinguishing prediction of Self-FEP Memory against
emotion-only models (CMR3) and against ad-hoc self-relevance heuristics
is the \emph{multiplicative} interaction between emotion and self
(E2). Bower's associative-network account of mood-congruent memory
\cite{bower1981mood} naturally predicts additivity: an emotional node
spreads activation, and a self node spreads activation, and the
contributions sum. Our model predicts the contributions \emph{combine}:
a vivid AND self-relevant memory has substantially lower forget\_score
than the additive prediction. The 0.71-magnitude interaction signature
recovered in E2 is the operational signature of this difference.

\paragraph{Clinical implications.}
E3a's $\piself{\to}0$ collapse of the SRE has empirical correlates.
Patients with dissociative identity disorder
\cite{conway2005smemsys} and severe depression show attenuated
self-reference effects in memory tasks. Our model expresses this as a
\emph{precision} deficit: when the self prior is unreliable, all
memories drift toward neutral self-relevance, and the simulated SRE
flattens. This gives a single computational lever --- $\piself$ --- that
\emph{can} express both ordinary SRE behaviour and attenuated SRE
behaviour within one framework. The clinical mapping remains a
prediction: impaired self-precision should reduce the self-reference
advantage in patient populations; this has not yet been fit to patient
data.

\paragraph{Self-evidencing and the memory system.}
The FEP-self literature \cite{apps2014freeself,limanowski2013minimalself}
emphasises that an agent samples the world to confirm its model of
itself (\textit{self-evidencing} \cite{hohwy2016selfevidencing}). Our
mechanism plugs into this naturally: encoding-time
$\muself$-updates only fire on user content that contains first-person
tokens, and the resulting precision is the agent's confidence that its
internal model of self correctly predicts what shows up about itself.
The retrieval-side $w_s\,\cos(e_i,\muself)$ term then biases recall
toward self-confirming content --- a state-dependent retrieval mechanism
in the encoding-specificity tradition \cite{tulving1973encodingspec}.

\paragraph{Where reconsolidation fits.}
Nader and Hardt \cite{nader2011reconsolidation} review reconsolidation:
retrieved memories become labile and re-encode under the current context.
Our retrieval loop already increments the retrieval count, but it does
not yet re-encode content. A natural extension is a
\texttt{re\_encode\_on\_recall} hook that blends the current $\muself$
and current emotional state into the recalled memory's stored fields,
with learning rate scaled by $\piself$. We flag this as v2 work in
\cref{sec:limitations}.

\paragraph{Position vs.\ CMR3.}
CMR3 \cite{cohen2022cmr3} models mood-congruent memory and intrusive
memory via context-emotion binding without a self representation; our
model adds the self dimension orthogonally. On self-referenced material
the two should diverge --- Self-FEP Memory predicts a robust effect,
CMR3 predicts none. On purely emotion-valenced material (no self
content) the two should agree closely. We have not yet performed this
quantitative comparison; \cref{sec:limitations} (L3) is candid about
this gap.
